\section{Discussion}
\label{sec:discussion}

\textbf{Language as the substrate.} Failure localization (C1), pre-execution review
(C2), and selective re-execution (C3) are not design choices about the Canvas App. They
are structural consequences: (1)~C1 follows from Local Addressability --- every step
has an exact flow index; its inputs are exactly its declared scope. (2)~C2 follows from
Interpretability + Enforceability --- the \texttt{.ncn} is human-readable, and what the
reviewer reads is exactly what executes. (3)~C3 follows from Local Addressability +
Enforceability --- the scope rule determines which downstream nodes depend on any step;
the orchestrator enforces this boundary. Existing LLM orchestration tools can add
checkpoint UI features; they cannot produce these properties (C1--C3) structurally
without first providing a reasoning medium that satisfies these six properties.

\textbf{Sustainable improvement through two-level CBR.} (1)~\textit{Level~1
accumulation}: every run deposits concrete cases; past checkpoints are available for
retrieval, reuse, and revision without operator intervention. (2)~\textit{Level~2
refinement}: each domain that warrants a new plan extends the abstract case library;
plans are revised as task patterns evolve and retained. (3)~\textit{Recursive
improvement}: the compilation pipeline --- itself a Level~2 case --- improves through
the same revision cycle, producing better plans, better Level~1 cases, and better
Level~2 revisions in turn. This distinguishes NormCode from a logging tool: it
accumulates experience in a form that can be retrieved, reused, and improved
\cite{bergmann2002,muller2015b}.

\textbf{What CBR gains from a reasoning medium.} The scope rule, determined at compile
time and enforced at runtime, guarantees that every checkpoint is a self-contained case.
The Level~1 case definition extends POCBR's case representation
\cite{leake2008,bottrighi2016} to resumable execution environments: a trace records
what happened; a suspended runtime is a state from which execution can be continued
with exactly the same data-dependency behavior, providing a guarantee that trace
retrieval does not. Finally, this paper concretely deploys the two-level structure
introduced by Muller \& Bergmann \cite{muller2015b}: four abstract reasoning patterns
(Level~2) each generating concrete cases on every run (Level~1), with the compilation
pipeline itself a Level~2 case.

\subsection{Limitations}

\textbf{Operational limitations.} Case retrieval currently requires operator judgment to
select the appropriate fork target; learned similarity measures would lower this barrier.
The case base grows without bound --- competence-based pruning is a natural next step.
Semantic checkpoints (LLM outputs) are bounded-stochastic: same input produces the same
output distribution, not a fixed value. Checkpoints whose tensors contain perceptual
signs are structurally self-contained, but reuse requires the referenced external
resources to be accessible; the case encodes \textit{what was referenced}, not its
content.

\textbf{Methodological limitations.} The Derivation phase uses an LLM to extract
conceptual structure from natural language, introducing non-determinism in plan
generation; a more formal derivation procedure is a direction for future work.

\textbf{Evaluation scope.} Comprehensive quality evaluation would require a
combinatorial test matrix ($|L_2| \times |L_1|$ cells), where outcome quality also
depends on LLM model, input domain, and subjective judgment. The evaluation in
Sect.~\ref{sec:evaluation} therefore argues for Level~2 generalizability --- that the
structural properties hold across the four production plans and that C1--C3 (failure
localization, pre-execution review, selective re-execution) are realized --- rather than
for absolute output quality, which varies with domain, model
choice, and subjective criteria.


